\documentclass[10pt]{article}
\usepackage[margin=20mm]{geometry}
\usepackage{amsmath,amssymb,mathtools,fontspec,unicode-math,xurl,hyperref}
\IfFontExistsTF{DejaVu Serif}{\setmainfont{DejaVu Serif}}{\setmainfont{Latin Modern Roman}}
\setmathfont{latinmodern-math.otf}
\setlength{\parindent}{0pt}\setlength{\parskip}{6pt}\sloppy
\title{Initial scientific review: M25-0044
The structure of isoperimetric bubbles}
\author{}
\date{25 October 2026}
\begin{document}
\maketitle
Abstract: One Minor issue is established: the reflection construction in Remark 9.14 does not in general preserve stability under the paper’s volume constraints. A two-disk example gives index form \(-\pi\) and an exactly volume-preserving path with negative second derivative. This invalidates the remark’s proposed necessity example, not Theorem 1.9. The inspected sphericity sign calculation, discrete maximum principle, dimension bound and complete-interface classification are consistent with their stated inputs. Coverage and unverified dependencies are explicit.\par
\section*{Identity, version and outcome}
Emanuel Milman and Joe Neeman, The structure of isoperimetric bubbles on \(\mathbb R^n\) and \(\mathbb S^n\), Acta Mathematica 234 (2025), 71–188, DOI `10.4310/ACTA.2025.v234.n1.a2`. The supplied publisher PDF, `papers/0044/0044.pdf`, is the reviewed version. Printed page equals one-based PDF page plus 70. Its SHA-256 is\par
> c1278e1d2fb399f769fec8f3f477421f20
> e27b9d0087328af561a532a5e708a6.\par
The hash above is line-broken for readability; the machine-readable source binding is authoritative. The entire supplied manifest was parsed and bound files verified before review. This is not counted as scientific inspection.\par
Outcome: one supported Minor finding, M01; zero confirmed Major findings. M01 concerns a peripheral example asserting necessity of an assumption. No damage to a proved main conclusion has been demonstrated. No whole-paper correctness certificate is issued.\par
\section*{Definitions, coverage and main dependencies}
A \(q\)-cluster includes the exterior cell; a single bubble therefore has \(q=2\). Individual cells need not be connected. In \(\mathbb R^n\) only the exterior cell has infinite volume. Total perimeter counts each interface once. Reflection symmetry means invariance of each labelled cell, not just invariance of the union of interfaces.\par
Definitions 2.26 and 2.28, printed p. 100 (PDF 30), require stationarity and \[\delta_X^1 V=0\quad\Longrightarrow\quad
 Q(X)=\delta_X^2 A-\langle\lambda,\delta_X^2V\rangle\geq0
 \qquad (X\in C_c^\infty).\] Perturbations are not restricted to reflection-symmetric fields or to preserving each connected component’s volume.\par
Theorem 1.3 concerns bounded stable regular clusters with reflection symmetry. The stable-cluster extension in Theorem 1.9 additionally assumes connected total boundary and finitely many cell components. The global minimizing conclusion in Theorem 1.13 is \(2\leq q\leq\min(5,n+1)\). Statements 1.14 and 1.16 are explicitly announced without proof in this paper; they are not treated as completed proofs here.\par
Text inspected: PDF pp. 1–18, 29–31, 39–41, 62–69, 83–91 and 97–110 (55 of 118 pages, including contents and figures). Actual source images inspected: PDF pp. 1, 8, 11, 12, 30, 64, 66, 67, 68, 83, 84, 87, 89, 91, 97, 98, 105, 106 and 108. Reading a page does not mean independently reconstructing every theorem it cites.\par
\section*{Problem 1: reflected copies need not be stable (M01)}
\subsection*{Source claim and exact scope}
Remark 9.14, printed p. 161 (PDF 91), proposes taking a stable \(q\)-cluster whose \(q-1\) finite cells lie in one open hemisphere and reflecting it about the equator. It states that the resulting \(q\)-cluster is stable and has disconnected cells, as an example showing why connectedness of the total boundary is needed in Theorem 1.9.\par
The suspect step is preservation of stability when a cell and its reflected copy receive the same label. The volume constraint then permits transferring volume between those two components.\par
\subsection*{An explicit input satisfying the construction}
Take \(M=\mathbb R^2\), \(n=2\), \(q=2\), equator \(\{y=0\}\), and \[D_+=B((0,2),1),\qquad D_-=B((0,-2),1).\] Start with the standard single bubble \((D_+,\mathbb R^2\setminus\overline{D_+})\), understood up to its null boundary. Its finite cell lies entirely in \(\{y>0\}\). It is bounded in the paper’s sense (compact total boundary), regular and stationary. It is stable: for a unit circle and a normal variation \(f\) of mean zero, \[Q(f)=\int_0^{2\pi}\bigl((f')^2-f^2\bigr)\,d\theta\geq0\] by the Fourier expansion, since every nonzero frequency has \(m^2-1\geq0\). Tangential fields do not change the constrained index.\par
Reflect and retain the same cell labels, exactly as required to obtain a \(q=2\) cluster with disconnected finite cell: \[\Omega_1=D_+\cup D_-,\qquad
 \Omega_2=\mathbb R^2\setminus(\overline{D_+}\cup\overline{D_-}).\] This is a regular stationary cluster; both circles have outward curvature 1 and the same pressure difference. It is bounded, each labelled cell is reflection invariant, and it has finitely many components. Also \(q=2\leq n+1=3\). It is not asserted to satisfy stability: that is the conclusion of the construction which the next calculation disproves.\par
\subsection*{Two exact checks of instability}
Equation reference: eq:disk. For \(|t|<1/2\) keep the centers fixed and set \[r_+(t)=\sqrt{1+t},\qquad r_-(t)=\sqrt{1-t}.\] The two disks stay disjoint and in their respective half-planes. Their total area is exactly \[V_1(t)=\pi r_+(t)^2+\pi r_-(t)^2=2\pi.\] The cluster perimeter is \[A(t)=2\pi\bigl(\sqrt{1+t}+\sqrt{1-t}\bigr),\qquad
 A'(0)=0,\qquad A''(0)=-\pi<0.                 \] This is an admissible smooth volume-preserving deformation near the compact interfaces, extendible by cutoffs in disjoint neighborhoods.\par
The same conclusion follows directly from the paper’s definition. Choose a smooth compactly supported ambient field with normal component \(f=1/2\) on the upper circle and \(f=-1/2\) on the lower one. Disjoint annular neighborhoods allow such a field, for example using local radial fields and cutoffs equal to 1 near each circle. Its first variation is \[\delta_XV_1=(1/2)(2\pi)-(1/2)(2\pi)=0,\qquad
 \delta_XV_2=-\delta_XV_1=0.\] There are no triple junctions, \(\operatorname{Ric}=0\), and \(|\mathrm{II}|^2=1\). Formula (2.15) gives \[Q(X)=\sum_{\pm}\int_{\partial D_\pm}(0-f^2)\,ds
      =-2(2\pi)(1/4)=-\pi.\] Thus the claimed stable output is unstable under the paper’s actual definition. Both computations are analytic; an independent exact-rational script also verifies the coefficient and index arithmetic.\par
\subsection*{Alternative readings and precise consequence}
Constraining the two disks’ volumes separately removes this variation, but changes the cluster’s labelled-volume constraints. Giving the reflected copy a different label also changes \(q\) and makes reflection exchange cells, contrary to the per-cell symmetry definition. Requiring only symmetric perturbations would define a weaker notion of stability than Definition 2.28. None of these readings validates the printed example under the existing definitions.\par
The example refutes the asserted stability-preservation construction, not Theorem 1.9 and not the existence of every possible different counterexample to dropping its connected-boundary assumption. Remark 9.14 follows the completed proof and is not used to establish the theorem or the minimizer classification. The demonstrated impact is therefore local and supports Minor. The precise safe correction is to remove the unsupported construction and leave necessity of that hypothesis unclaimed pending an example satisfying the actual stability condition. This review does not propose removing the hypothesis from the theorem.\par
The inspected arXiv v4 (20 April 2025) retains the same Remark 9.14. No correction of this point was identified in the inspected author publication listing; this limited search is not proof of novelty or absence of all errata.\par
\section*{Checks of the central arguments}
\subsection*{The sign producing perpendicular sphericity}
On printed pp. 135–137 (PDF 65–67), put \(x=\langle N,p\rangle\) and \(z=\langle N,c\rangle\), with \(c=\nu-kp\). Lemma 7.8 gives \[L_{\rm Jac}x=|\mathrm{II}_0|^2x-(n-1)kz.\] Trace stability yields \(0\leq-\int x^2|\mathrm{II}_0|^2+(n-1)\int kxz\). The cutoff conformal-field estimate on the sphere yields \(0\leq-(n-1)\int\langle N,\nu\rangle z\). Since \(\langle N,\nu\rangle=kx+z\), adding cancels the mixed term: \[0\leq-\int\bigl(x^2|\mathrm{II}_0|^2+(n-1)z^2\bigr)\leq0.\] The sign is correct. Stationarity bounds the finitely many \(k_{ij}\), so \(|c_{ij}|^2=1+k_{ij}^2\) controls the cutoff error. In Euclidean space the required zero integral comes from Lemma 7.9, whose proof explicitly uses global \(L^2\) curvature control, rather than merely control away from the singular set. Full proofs of the earlier approximation and integrability results remain dependencies.\par
There is an independent check of the Euclidean isotropy formula in Remark 7.10. For a trace-free symmetric matrix \(B\), the flow \(\exp(tB)\) preserves every cell volume. With a cutoff agreeing near the compact boundary, stationarity and the surface Jacobian give \[0=\int_{\Sigma^1}\bigl(\operatorname{tr}B-\nu^TB\nu\bigr)\,dA
   =-\operatorname{tr}\left(B\int_{\Sigma^1}\nu\otimes\nu\,dA\right).\] The integral is therefore a scalar multiple of the identity; its trace is \(A(\Omega)\), giving \(A(\Omega)I/n\). This checks the normalization without assuming isotropy from a picture.\par
\subsection*{Connectedness and the dimension restriction}
The strict maximum statement of Lemma 9.12, p. 157 (PDF 87), uses positive edge weights and connectedness after deleting the source vertex. On the set \(W\) of maximum values of the solution of \(L_Aa=e_1-e_V\), the sink cannot occur. Positivity of the boundary flux forces \(1\in W\), and every vertex of \(\partial W\) must be 1. If \(W\) contained any other vertex, deleting 1 would disconnect the graph. Hence \(W=\{1\}\), establishing the strict inequality actually used.\par
The cusp argument on pp. 158–159 uses \(h_{ij}=a_{ij}|\langle N,p\rangle|\) with zero first volume variation and \(Q(h)=0\). Proposition 9.2 compares quadratic integrands almost everywhere; it does not apply a smooth integration-by-parts formula directly to the cusp. Elliptic regularity then contradicts its nonzero normal derivative jump. This reasoning is conditional on the scalar-field approximation and regularity lemmas, not on a claim that the cusp is initially smooth.\par
For Proposition 9.13, \(s\leq q-1\leq n\) equatorial cells leave a nonzero direction orthogonal to the first \(s-1\) constraint normals; its sign can be chosen to satisfy the last inequality. A ray in that direction stays in the open polyhedron and reaches the unit sphere, contradicting density of the equatorial cells. The argument explains the bound \(q\leq n+1\) and does not extend automatically to \(q=n+2\).\par
\subsection*{Complete interfaces and small adjacency graphs}
For Proposition 10.10, pp. 167–168 (PDF 97–98), all interfaces imply \(G(e_i-e_j,e_i-e_j)=1\). To verify that these data determine \(G\), set \(u_i=e_i-e_q\). A symmetric form \(B\) vanishing on all \(e_i-e_j\) satisfies \(B(u_i,u_i)=0\) and \(B(u_i-u_j,u_i-u_j)=-2B(u_i,u_j)=0\), hence \(B=0\). Consequently \(G=\tfrac12 I\) on the sum-zero subspace, with the paper’s factor \(\tfrac12\) correct.\par
Corollary 11.3 uses Taylor’s classification only when the effective cone dimension is at most 3. In the remaining \(d=4\), \(q=5\) case, affine independence of five generators gives a simplex normal fan, so every pair shares an interface; it does not presume classification of all four-dimensional minimizing cones.\par
The graph list in Figure 11.5 is complete: a five-vertex graph of minimum degree 3 has complement of maximum degree 1, hence a matching of size 0, 1 or 2. Independent enumeration gives respectively 1, 10 and 15 labelled graphs, exactly three isomorphism types. The subsequent geometric exclusions use minimization and contact regularity; the graph count alone is not their proof. The angle inequality in Lemma 11.6 gives \(120(q-1)<180(q-3)\), hence \(q>7\), consistent with the stated nondegenerate range.\par
\section*{Later clarification, limitations and conclusion}
The authors’ subsequent paper [boundary], Remark 1.6, explicitly identifies the simple-connectedness input implicitly used in the volume-preserving boundary-motion argument of Lemma 9.3, printed pp. 153–154. The original model spaces in dimensions \(n\geq2\) have that property. This is recorded as a known clarification, separately from M01; no counterexample to Lemma 9.3 in its stated model spaces is established. Only the later introduction and stated topological input were inspected, not its full proof.\par
Unperformed work includes the full regularity and scalar-to-vector approximation machinery, the complete simplicial-homology/Voronoi construction in Section 8, all singular maximum-principle details, the omitted properness and differential arguments in the volume-parameterization sketch, and all external foundations. No author code was supplied in the manifest or executed. No empirical experiment is relevant to the explicit counterexample.\par
Judgment | One Minor: false stability claim in the reflected-copy example. No established main-theorem failure.\par
Repair | Deletion of that unsupported example is safe; no replacement necessity theorem verified.\par
Adoption | No author contact or adoption request. Human scientific acceptance pending.\par
Remaining | Necessity of the connected-boundary hypothesis requires separate evidence; unverified dependencies above remain open in this review.\par
\section*{Supporting checks}
Reviewer-written code is saved in the evidence directory as `independent\_checks.py`. It uses only Python standard-library exact fractions and enumeration of the \(2^{10}\) labelled five-vertex graphs. It was actually executed by this conversation with existing Windows Python; output is `independent\_checks\_result.json`, SHA-256\par
> 7fa9ba2cb15031cb2e52ffe6bf7888eb2230
> 205e0dbf66afb8a48d923f2b1a52.\par
It reports \(A''(0)/\pi=-1\), \(Q/\pi=-1\) and graph counts \(1,10,15\). The derivation, geometry and assumption checks are given above; the finite computation alone would not prove them. Source text, displayed-page PNGs, read/render operation records and primary-source lookups are retained in the same evidence directory. Extraction and index parsing are not review.\par
No author code was executed and no compiler executable was installed.\par
Human scientific adjudication remains pending.\par
9 E. Milman and J. Neeman, The structure of isoperimetric bubbles on \(\mathbb R^n\) and \(\mathbb S^n\), Acta Math. 234 (2025), 71–188. DOI: `10.4310/ACTA.2025.v234.n1.a2`. Assigned publisher PDF.\par
E. Milman and J. Neeman, same title, arXiv:2205.09102v4, 20 April 2025. https://arxiv.org/abs/2205.09102v4. Version metadata and Remark 9.14 compared on 25 October 2026.\par
E. Milman and J. Neeman, On the connectedness of a minimizing cluster’s boundary, arXiv:2512.15403v1, 17 December 2025. https://arxiv.org/abs/2512.15403v1. Introduction, particularly Remark 1.6; consulted 25 October 2026.\par
E. Milman and J. Neeman, Plateau Bubbles and the Quintuple Bubble Theorem on \(\mathbb S^n\), arXiv:2307.08164v3, 28 December 2024. Metadata and abstract only; proof not reviewed.\par
\end{document}
